\newcommand{\figuraMoleculasVolumen}{%
\begin{figure}[htbp]
    \centering
    \begin{tikzpicture}[>=Latex,font=\small]
        % Recipiente grande y volumen de observación.
        \draw[very thick,rounded corners=3pt,fill=blue!3]
            (0,0) rectangle (6.2,4.2);
        \node[font=\bfseries] at (3.1,4.55)
            {Gas con $N_{\mathrm{tot}}$ moléculas};

        \foreach \x/\y in {
            0.35/0.45,0.85/1.25,0.55/2.20,1.15/3.25,1.65/0.65,
            1.85/1.65,1.55/2.75,2.25/3.65,2.55/0.35,2.85/0.95,
            2.55/1.55,3.15/1.95,3.65/2.65,3.85/1.25,4.35/0.55,
            4.65/1.55,4.25/2.25,4.85/3.35,5.35/0.85,5.65/1.95,
            5.35/2.75,5.85/3.65,3.25/3.55,0.75/3.75}
        {
            \fill[blue!65] (\x,\y) circle (0.075);
        }

        \draw[very thick,red!70!black,fill=red!5,rounded corners=2pt]
            (2.20,1.05) rectangle (4.15,3.15);
        \node[font=\bfseries,red!70!black] at (3.175,3.48)
            {volumen $V$};

        \foreach \x/\y in {
            2.55/1.55,3.15/1.95,3.65/2.65,3.85/1.25,3.25/3.00}
        {
            \fill[red!75] (\x,\y) circle (0.095);
        }
        \node[align=center] at (3.175,0.62)
            {$m$ moléculas observadas\\$p=V/V_{\mathrm{tot}}$};

        \draw[->,very thick,red!70!black]
            (6.55,2.10)--node[above]{muchas repeticiones}(7.75,2.10);

        % Distribución de Poisson para lambda=3.
        \begin{scope}[shift={(8.15,0.35)}]
            \draw[->,thick] (0,0)--(5.4,0) node[right]{$m$};
            \draw[->,thick] (0,0)--(0,3.65) node[left]{$P(m)$};

            \foreach \m/\x/\h in {
                0/0.35/0.65,1/0.85/1.94,2/1.35/2.91,3/1.85/2.91,
                4/2.35/2.18,5/2.85/1.31,6/3.35/0.66,7/3.85/0.28,
                8/4.35/0.11,9/4.85/0.04}
            {
                \fill[blue!30] (\x-0.18,0) rectangle (\x+0.18,\h);
                \draw[blue!70!black,thick]
                    (\x-0.18,0) rectangle (\x+0.18,\h);
                \node[font=\scriptsize,anchor=north] at (\x,-0.04) {$\m$};
            }

            \draw[dashed,red!75!black,thick]
                (1.85,0)--(1.85,3.35);
            \node[red!70!black,anchor=south] at (1.85,3.35)
                {$\langle m\rangle=\lambda=3$};
            \draw[<->,very thick,orange!85!black]
                (0.98,0.38)--node[above]{$\Delta m=\sqrt{3}$}(2.72,0.38);
            \node[font=\bfseries] at (2.70,4.12)
                {Distribución de Poisson};
        \end{scope}

        \node at (6.7,-0.65)
            {$\displaystyle
             P(m)=\frac{\lambda^m e^{-\lambda}}{m!}
             \qquad
             \frac{\Delta m}{\langle m\rangle}=\frac{1}{\sqrt{\lambda}}$};
    \end{tikzpicture}
    \caption{Conteo de moléculas en un volumen pequeño. Cada molécula del gas
    tiene una probabilidad $p$ de encontrarse dentro de $V$. Al repetir la
    observación, el número $m$ fluctúa alrededor de $\lambda=N_{\mathrm{tot}}p$
    y sigue, en el límite diluido, una distribución de Poisson.}
    \label{fig:taller-poisson-volumen}
\end{figure}%
}
